\documentclass[11pt]{article}
\usepackage[margin=1in]{geometry}
\usepackage{amsmath,amssymb,amsthm}
\usepackage{booktabs}
\usepackage{array}
\usepackage{enumitem}
\usepackage{xcolor}
\usepackage{hyperref}
\hypersetup{colorlinks=true,linkcolor=blue!50!black,urlcolor=blue!50!black,citecolor=blue!50!black}
\usepackage{fancyhdr}
\usepackage{listings}
\usepackage[protrusion=true,expansion=false]{microtype}

\lstset{
  basicstyle=\ttfamily\footnotesize,
  breaklines=true,
  frame=single,
  framesep=4pt,
  backgroundcolor=\color{gray!6},
  columns=fullflexible,
  keepspaces=true,
}

\pagestyle{fancy}
\fancyhf{}
\fancyhead[L]{\small Section~9 cone chain \& $4/3$ kill tests}
\fancyhead[R]{\small TR-2026-FF06-I7}
\fancyfoot[C]{\thepage}

\newtheorem{definition}{Definition}
\newtheorem{remark}{Remark}

\title{\textbf{Finite-Model Kill Tests for a Gap--Cone Dependency Chain\\
and the $4/3$ Coincidence}\\[0.35em]
\large Diagnostic note accompanying ACS Framework Issue~\texttt{\#7}}
\author{B.\ Wallace \quad (TensorRent)\\ \small with computational verification in-repo}
\date{\small TR-2026-FF06-I7 \ ---\ 2026-07-22 \ ---\ SIP License v1.1}

\begin{document}
\maketitle

\begin{abstract}
Two diagnostic programmes are reported as finite-model kill tests, not
as theorems.  The first asks whether a proposed dependency direction
\[
  \text{spectral gap}\ \to\ \text{clustering}\ \to\ \text{cone-sharpness proxy}
\]
is supported as a robust monotone relation on small lattice models
(TFIM, XXZ+$h$, long-range Ising) and on an exact float-free classical
Ising analogue.  The second asks whether the numeric coincidence between
an ACS-side constant $\beta=4/3$ and a density-side exponent form
$\alpha(d)=1+1/d$ at $d=3$ survives mechanism-level robustness checks.
Across the tested windows the Section~9 chain is not supported as a
robust monotone relation; the $4/3$ equality fails normalisation
robustness and collapses to the identity family $n=d+1$ on the scanned
integer lattice.  All scripts, JSON artefacts, SHA-256 anchors, and an
end-to-end verification pipeline are included in-repo.  Claims are
scoped to tested models and parameter windows.
\end{abstract}

\section*{License}
Co-governed and enforced under the
Sovereign Integrity Protocol License (SIP License v1.1):
\url{https://github.com/tensorrent/ACS-Framework/blob/main/LICENSE}.

\section{Scope and claim discipline}

This note reports finite tested outcomes only.  It does not claim a
universal no-go theorem, a proof of cone rigidity, or a derivation of a
critical coupling.  All findings are bounded by the tested model
families, parameter windows, and (where applicable) numeric backend.

\begin{remark}[Relation to Little-type programmes]
A dependency statement of the form ``static reduction is meaningful only
after observables are rigid enough to support a sharp cone'' motivates
the Section~9 toy, but the present work does not establish that
dependency as a theorem.  The tests ask whether a computable proxy of
the proposed direction survives kill criteria on finite models.
\end{remark}

\section{Research questions}

\paragraph{RQ-1 (Section~9 toy chain).}
In finite lattice models, does larger spectral gap co-vary with both
shorter static clustering length and lower outside-cone commutator
leakage (or an exact influence analogue)?

\paragraph{RQ-2 ($4/3$ mechanism).}
Is the equality at $4/3$ generated by a robust mechanism, or does it
collapse under normalisation changes and bookkeeping-identity checks?

\section{Methods}

\subsection{Section~9 float numerics}

Scripts under \texttt{code/issue7/}:
\begin{itemize}[leftmargin=1.4em,itemsep=2pt]
  \item \texttt{section9\_toy\_kill\_test.py} --- baseline TFIM toy;
  \item \texttt{exhaustive\_issue7.py} --- multi-size / multi-threshold sweep;
  \item \texttt{cross\_model\_kill\_pass.py} --- TFIM and XXZ+$h$ with dual leakage summaries;
  \item \texttt{long\_range\_kill\_pass.py} --- power-law Ising, $\alpha\in\{3,2\}$.
\end{itemize}

Observables: gap $\Delta=E_1-E_0$; clustering length $\xi$ from connected
$Z_0$--$Z_r$ correlators; outside-cone leakage from
$\|[Z_0(t),Z_r]\|$ (max and, where stated, mean).

Decision rule (float suite): support requires
$\mathrm{corr}(\mathrm{gap},\xi)<-0.2$ and
$\mathrm{corr}(\mathrm{gap},\mathrm{leak\_proxy})<-0.2$
for every required leakage proxy.

\subsection{Section~9 exact (no IEEE float)}

\texttt{section9\_exact\_kill\_test.py} uses integer Ising energies,
integer Boltzmann base $B$, connected correlators in
\texttt{Fraction}, and local-majority influence fractions in
\texttt{Fraction}.  The decision path uses only signs of
\texttt{Fraction} covariances.

Companion exact ACS TE automaton (already in-repo):
\texttt{code/acs\_codebase/extras/integer\_acs.py}.

\subsection{$4/3$ mechanism tests}

\texttt{mechanism\_4over3\_test.py} and the $4/3$ half of
\texttt{exhaustive\_issue7.py}:
\begin{enumerate}[leftmargin=1.4em,itemsep=2pt]
  \item K1 --- normalisation robustness: scale $\beta=4/3$ by $c$ and
        inspect $\hat d=1/(c\beta-1)$;
  \item K2 --- integer lattice scan of $\alpha(d)=1+1/d$ against
        $\beta(n)=n/(n-1)$.
\end{enumerate}

\subsection{End-to-end verification}

\texttt{verify\_issue7\_pipeline.py} reruns all Issue~\texttt{\#7}
scripts, rewrites canonical stdout logs, checks stdout-vs-JSON
consistency, and verifies SHA-256 anchors listed in the companion
markdown artefacts under \texttt{docs/issue7/}.

\section{Results}

\subsection{Section~9 baseline and exhaustive}

Baseline TFIM ($N=7$):
$\mathrm{corr}(\mathrm{gap},\xi)=-0.6375$,
$\mathrm{corr}(\mathrm{gap},\mathrm{leak})=+0.9636$
--- chain not supported.

Exhaustive float sweep: support rate $0/12$.
Gap--$\xi$ correlations remain negative; gap--outside-leak correlations
remain positive in every tested cell.

\subsection{Cross-model and long-range}

\begin{center}
\begin{tabular}{lcc}
\toprule
Family & Support & Rate \\
\midrule
TFIM (cross-model grid) & $0/6$ & $0.00$ \\
XXZ+$h$ & $1/6$ & $0.17$ \\
Long-range $\alpha=3$ & $2/6$ & $0.33$ \\
Long-range $\alpha=2$ & $3/6$ & $0.50$ \\
\bottomrule
\end{tabular}
\end{center}

The single XXZ support cell and the long-range support cells concentrate
at the smallest leakage threshold ($10^{-3}$).  Partial support is
threshold-dependent, not a uniform monotone trend.  Long-range Ising
--- where strict Lieb--Robinson cones are known to weaken --- shows the
highest partial-support fraction, consistent with threshold sensitivity
of the cone proxy rather than establishment of the chain.

\subsection{Exact float-free Section~9}

\begin{itemize}[leftmargin=1.4em,itemsep=2pt]
  \item $\mathrm{sign}\,\mathrm{cov}(\mathrm{gap},\mathrm{cluster\_weight})=-1$;
  \item $\mathrm{sign}\,\mathrm{cov}(\mathrm{gap},\mathrm{outside\_affect})=0$
        (outside-cone affect identically zero under local majority);
  \item joint chain support: false.
\end{itemize}

\texttt{integer\_acs.py} on $\mathbb{Z}_{16}$: uncoupled/symmetric
$\Delta I\approx 0$; asymmetric $\Delta I\approx\pm 1.499$ with sign
reversal under $f\leftrightarrow g$.  Probabilities are exact rational.

\subsection{$4/3$ mechanism}

Baseline $\beta=4/3$ implies $\hat d=3$.  Under normalisation scaling,
inferred-dimension spread is $74$ on the tested grid, with non-finite
values on part of the interval.  Exact lattice hits all satisfy
$n=d+1$.  Mechanism not established on this surface.

\subsection{Verification}

End-to-end pipeline status: \texttt{PASS\_WITH\_CAUTION};
all scripts exit successfully; consistency errors $0$; documentation
hash checks $0$.

\section{Reproducibility}

Canonical artefacts and SHA-256 anchors are listed in
\texttt{docs/issue7/issue7\_paper\_section9\_and\_4over3.md}
and \texttt{docs/issue7/issue7\_takeup\_section9\_and\_4over3.md}.

Observed environment for float suite and anchors:
Python~$3.14.6$, NumPy~$2.5.0$,
\texttt{macOS-15.3.2-arm64-arm-64bit-Mach-O}.

\begin{lstlisting}[caption={Reproduction commands}]
python3 code/issue7/section9_toy_kill_test.py
python3 code/issue7/mechanism_4over3_test.py
python3 code/issue7/exhaustive_issue7.py
python3 code/issue7/cross_model_kill_pass.py
python3 code/issue7/long_range_kill_pass.py
python3 code/issue7/section9_exact_kill_test.py
python3 code/issue7/verify_issue7_pipeline.py
\end{lstlisting}

\section{Limitations}

\begin{itemize}[leftmargin=1.4em,itemsep=2pt]
  \item Float NumPy Section~9 scripts remain single-platform
        anchored.\footnote{Float NumPy Section~9 scripts (exact
        diagonalisation / commutator norms) are anchored to one machine
        (macOS~$15.3.2$ arm64, Python~$3.14.6$, NumPy~$2.5.0$), with
        same-machine byte-identical reruns recorded in
        \texttt{docs/issue7/issue7\_verification\_report.json}.
        The exact Section~9 kill test and
        \texttt{integer\_acs.py} use integer/\texttt{Fraction} on the
        decision path; for those diagnostics, cross-machine IEEE float
        drift is not a concern and a second platform is not required.
        The $4/3$ lattice/identity checks are likewise exact.}
  \item Exact Section~9 / integer ACS diagnostics are float-free on the
        decision path.
  \item Finite-size and finite-window coverage can miss alternate regimes.
  \item No claim is made beyond tested models and tested parameter windows.
\end{itemize}

\section{Conclusion}

Within the reported finite tests, the Section~9 dependency chain is not
supported as a robust monotone relation across the tested model families
and thresholds.  The long-range family shows the highest partial-support
fraction ($5/12$), concentrated at the smallest leakage threshold --- a
threshold-dependence effect, not evidence of a stable universal trend.
The exact float-free analogue also fails the joint chain criterion
(cluster covariance sign correct; outside-cone affect identically zero
under local majority).  The $4/3$ equality does not satisfy the tested
mechanism robustness criteria and is consistent with an identity-family
explanation on the scanned lattice.  The package is reproducible
in-repo with anchored artefacts; end-to-end verification retains caution
only for the remaining float NumPy Section~9 scripts.

\section*{Tier map (MANIFEST)}

\begin{center}
\begin{tabular}{lll}
\toprule
Claim & Tier & Evidence \\
\midrule
TFIM exhaustive chain support rate $0/12$ & T1 & \texttt{exhaustive\_issue7.py} \\
Cross-model combined support $1/12$ & T1 & \texttt{cross\_model\_kill\_pass.py} \\
Long-range combined support $5/12$ & T1 & \texttt{long\_range\_kill\_pass.py} \\
Exact joint chain support false & T1 & \texttt{section9\_exact\_kill\_test.py} \\
$4/3$ identity family $n=d+1$ & T1 & \texttt{mechanism\_4over3\_test.py} \\
$4/3$ mechanism established & T4 & failed K1/K2 \\
Section~9 chain as general theorem & --- & not claimed \\
\bottomrule
\end{tabular}
\end{center}

\end{document}
